% Requires \usepackage{booktabs}. Full-width floats for two-column conference layouts.
% Source tables: S1
\begin{table*}[t]
\centering
\small
\setlength{\tabcolsep}{4pt}
\caption{Held-out test comparison of the current FIA-off/on pair.}
\label{tab:selected-s1}
\begin{tabular}{lrrrrrr}
\toprule
Model & AP & AP50 & AP75 & APs & APm & APl \\
\midrule
G+B & 24.111 & 41.668 & 24.197 & 14.334 & 34.136 & 44.423 \\
Full & 24.229 & 42.009 & 24.208 & 14.618 & 34.290 & 44.873 \\
Delta (pp) & 0.118 & 0.341 & 0.011 & 0.284 & 0.155 & 0.451 \\
\bottomrule
\end{tabular}
\par\smallskip
\begin{minipage}{0.97\textwidth}\footnotesize The checkpoints selected by native validation AP are evaluated once under the common test protocol. Test AP improves by 0.118 percentage points, compared with 0.832 points in native validation. These single-seed observations do not establish statistical significance.\end{minipage}
\par\medskip
\end{table*}

% Source tables: S2
\begin{table*}[t]
\centering
\small
\setlength{\tabcolsep}{4pt}
\caption{Best and final checkpoints under the common evaluator.}
\label{tab:selected-s2}
\begin{tabular}{lrrr}
\toprule
Model & Checkpoint & Val AP & Test AP \\
\midrule
D-FINE-L & best & 28.796 & 22.679 \\
D-FINE-L & final & 28.501 & 22.554 \\
Full & best & 30.727 & 24.229 \\
Full & final & 30.528 & 24.133 \\
\bottomrule
\end{tabular}
\par\smallskip
\begin{minipage}{0.97\textwidth}\footnotesize Best checkpoints were selected on native validation AP, whereas final denotes epoch 100. Test results were not used for checkpoint selection. Common-evaluator validation differs from the native training metric.\end{minipage}
\par\medskip
\end{table*}
